\section{Related Work}
\label{sec:related_work}

Two lines of work motivate our approach: interventional methods for diagnosing judge bias, and measurement theory on the sensitivity of different elicitation formats.

\paragraph{Interventional Approaches to Judge Bias}
Text-based interventional methods have a growing methodological foundation \citep{feder2021causalm,wu2021polyjuice,feder2022causal}.
\citet{reber2025rate} propose RATE, a framework that rewrites model outputs along style attributes using imperfect counterfactuals and measures the causal effect on reward model scores via Likert-scale Average Treatment Effect (ATE).
RATE applies a correction term to account for rewriting imperfections, but its reliance on Likert ATE as the primary metric is a vulnerability: our experiments show that Likert scoring exhibits a ceiling effect in which 60\% of style-transformed pairs receive identical integer scores, yielding non-significant ATE estimates ($p = 0.332$) for the same interventions that pairwise comparison detects at $p = 0.006$.
Whether this format sensitivity extends to RATE's broader attribute space is an open question; our results motivate independent validation of Likert-based effect estimates alongside pairwise measures.
CSD differs from RATE in three respects: it targets individual style axes rather than general attributes, it measures bias through position-controlled pairwise comparison rather than Likert ATE, and it enforces content preservation via NLI entailment gating rather than post-hoc correction for imperfect rewrites.
\citet{counterfactual2025rhetorical} also adopt an interventional LLM-based design, rewriting texts to measure rhetorical style effects and aggregating preferences with a Bradley-Terry model.
Their framework operates at persona-level granularity (e.g., ``writes like a professor''), whereas CSD intervenes on specific style axes (e.g., formality, conciseness) and uses a sign test with position control rather than parametric aggregation.

\paragraph{Measurement Sensitivity in NLG Evaluation}
The distinction between ordinal and interval measurement dates to \citet{stevens1946theory}, whose taxonomy of scales remains foundational.
The superiority of pairwise comparison over rating scales for eliciting reliable human judgments has been established since \citet{thurstone1927law}, whose Law of Comparative Judgment formalizes why forced-choice comparisons yield finer discriminability than absolute ratings; \citet{carterette2008here} confirm this for relevance assessment in information retrieval.
\citet{preston2000scale} and \citet{simms2019does} show that scale granularity affects reliability and psychometric precision; item response theory provides a complementary framework for quantifying measurement precision as a function of item characteristics \citep{embretson2000irt}.
\citet{krosnick1999survey} documents satisficing behavior that may compound the tie inflation we observe.
\citet{decoster2011median} show that artificial categorization of continuous variables discards variance and reduces statistical power, directly relevant to the Likert Blindspot mechanism.
Within NLP, \citet{amidei2019likert} review the use of rating and Likert scales \citep{likert1932technique} in NLG evaluation, identifying inconsistencies in scale design and reporting that compromise cross-study comparability.
\citet{kiritchenko2017bws} and \citet{novikova2018rankme} confirm that forced-choice methods detect quality differences that rating scales miss; \citet{deutsch-etal-2023-ties} show that tie handling directly affects metric sensitivity, consistent with our tie-inflation finding.
Our contribution is the first controlled quantification of this per-axis sensitivity gap in the LLM judge setting, identifying the two-layer mechanism (tie inflation plus residual signal attenuation) that drives it.

\paragraph{Style Bias Detection and Mitigation}
The foundational LLM-as-judge framework was established by \citet{zheng2023judging}, who introduced MT-Bench and Chatbot Arena and documented position bias in pairwise evaluation.
\citet{wang2023faireval} further showed that preference rankings can be manipulated by swapping option order, proposing a win/tie/lose annotation scheme, and \citet{shi2024judging} systematically study position bias across 15 judges, showing it varies by model and quality gap.
\citet{li2024calibraeval} address selection bias through calibration of prediction distributions during inference.
These works focus on \emph{position} bias and calibration; \citet{ye2024justice} systematically quantify 12 bias types using automated counterfactual modifications, and \citet{koo2024cobbler} benchmark cognitive biases in LLM evaluators, though neither provides per-axis isolation of individual style effects.
Style bias as a separate, dominant bias source has received attention only recently \citep{wu2023style,murugadoss2025sos,li2026preference,soumik2026judging}; \citet{saito2023verbosity} specifically quantify verbosity bias in multi-turn chatbot evaluation, and \citet{chen2024humans} compare human and LLM judge biases, finding partially overlapping but distinct bias profiles; \citet{bavaresco2025llms} show that LLM--human agreement varies substantially across tasks.
\citet{li2026preference} further show that style-mediated \emph{preference leakage} propagates through the evaluation pipeline, motivating per-axis isolation of style effects.
\citet{tan2025judgebench} provide JudgeBench, a benchmark revealing that even strong judges perform near chance on challenging pairs, underscoring the need for controlled experimental designs rather than aggregate accuracy metrics.
Multi-feature frameworks such as CASSE \citep{casse2026concise} and FairJudge \citep{wang2026fairjudge} decompose style into finer dimensions but rely on correlational analyses---they cannot distinguish whether observed associations reflect causal style effects or confounds with content quality, which our NLI-gated rewriting design addresses.
\citet{stureborg2024inconsistent} demonstrate inconsistency and bias across evaluation formats, and \citet{park2024offsetbias} show that verbosity and other surface features systematically inflate judge scores.
Crucially, \citet{jeong2025comparative} identify the Comparative Trap: pairwise comparison amplifies style biases relative to pointwise scoring, raising the concern that our pairwise-based detection may partly reflect format-level amplification rather than pure detection.
We address this directly in \S\ref{sec:amplification}.

\paragraph{Bradley-Terry Models in Evaluation}
The Bradley-Terry model \citep{bradley1952rank} underpins Chatbot Arena \citep{zheng2023judging,chiang2024chatbot} and AlpacaFarm \citep{dubois2023alpacafarm}; we adopt BT because it uses all comparison data including inconsistent pairs, yielding proper effect-size CIs, and combine it with GEE decomposition to expose per-axis confound structures invisible to aggregate scores.

\paragraph{Mode-specific biases in LLM judges.}
Unlike human evaluators, LLM judges acquire format-specific biases through RLHF training: pairwise comparison serves as the reward signal in DPO and RLHF pipelines, while Likert scoring is primarily used for benchmark evaluation.
This asymmetric training exposure motivates independent quantification of the Likert--pairwise gap, as format-specific biases in reward models directly affect learned policies through the training objective.
